\documentclass[12pt]{article}
\usepackage{graphicx} % Required for inserting images
\usepackage{xcolor}
\usepackage[left=1in,top=1in,right=1in,bottom=1in]{geometry}

\begin{document}

\section{Table 1 on summary statistics}
\section{Tables 2 and 3}
My preference for the new Table 2 would be that we copy-paste the current Table 3, except we drop column (2). In effect, we replicate deHaan et al. (2025)'s Table 4, Panel A in the first two columns. Then, we show that those left behind experience significant outflows.

My preference for the new Table 3 would be that we copy-paste the current Table 2.

By presenting statistically significant results first, and non-results second, my hope is that the paper won't be perceived as a non-results paper.
\section{Table 4 on fee difference btwn livestreamed and not}
Can we do something similar to what we do in Table 10 for fund returns? Instead of ``Return'', ``Alpha (FF3)'', and ``Alpha (FF5)'', here we would report ``Management Fee'', ``Service Fee'', and ``Loads''. Hopefully, we will \textit{not} find (economically) significant difference between livestreamed and not in terms of fee components.
\section{Table 5}
My preference would be that we drop the current Tables 4--5, making the current Table 6 the new Table 5.

Perhaps, we can make the current Table 5 the first panel of the new Table 5, and make the current Table 6 the second panel of the new Table 5.
\section{Table 6 on \$ fee revenue}
We could present results on \$ fee revenue difference between livestreamed and not, similar to what we do in the new Table 4 for fund fees, but my worry is that such results may feel too repetitive. Thus, my preference would be to make the current Table 8 the new Table 6, and leave it at that.
\section{Tables 7 and 8}
Given that I'm suggesting we present fund-level results first and family-level results second, my preference would be to make the current Table 10 the new Table 7; the current Table 9 the new Table 8.
\section{Table 9}
The current Table 11 would be the new Table 9.

\section*{Rich's Comments}

Hey guys,

Ilias and I talked to Rich Evans earlier today. Good news is he's really excited about this data. Bad news is, basically, he wants to start a new project with his idea.

In particular, he wants us to explore learning by fund families (in contrast to investor learning, which is far more typical in the literature). Given that social media livestreaming is a novel form of marketing, he wants to explore how families change their livestream contents over time.

First, he asks: what are the key determinants of fund flows (e.g., fund performance, fund characteristics) before 2020, i.e., before the era of social media livestreaming.
He wonders whether fund families try to advertise funds based on these traditionally key determinants at the start of 2020, i.e., at the start of the social media livestreaming era.

He then asks: what kind of livestreams---what you talk about and how you talk about---attract fund flows at high frequency.
He wonders whether fund families learn over time what kind of livestreams actually work, i.e., actually attract fund flows. Here, he thinks both intraday share-class-level flows and Danmaku would be super useful. Basically, he wants us to look at what kind of livestreams actually work and whether fund families increase the frequency of such livestreams in the future. We told him that Danmaku is subject to a selection bias, but he thinks that's OK. Even analyzing for which livestreams fund families choose to allow Danmaku or not would be super interesting in itself.

Ultimately, he thinks it would be super interesting if we can show that the primary determinants of fund flows change over time because (i) the industry has shifted away from traditional forms of marketing to social media livestreaming, (ii) fund families learn how to optimize livestreaming for the purpose of attracting flows, and (iii) livestreaming brings in different types of investors than those that traditionally invested in mutual funds. He hypothesizes that all of this may have changed the primary determinants of fund flows.

Methodologically, basically, he wants us to do content analysis. What do fund families talk about in these livestreams (macro? economy? funds? etc.) and how do they talk about these things. This is closer to what others are doing, but we agree with him that it's very different in the sense that this would be the only paper that focuses on time variation (and convergence) in contents of social media livestreams.

To summarize, he wants us to empirically analyze whether there is "practice makes perfect," i.e., whether fund families learn how to more effectively market their funds to attract flows) such that there are some systematic shifts in livestream contents.

I think the common theme based on our chats with both Vikas and Rich is that, unfortunately, quarterly flows are too low-frequency. In turn, intraday share-class-level flow data or Danmaku, even if they are subject to selection bias and what not, is still preferable. In particular, Rich says that people analyze fund manager newsletters in the US and their effects on fund flows, which is similar in spirit to what we are doing with quarterly flows. What he says is rare about the data in China is the fact that the data on intraday share-class-level flow and Danmaku (even if they are subject to selection bias) provide us high-frequency observations of investor reactions. Such observations are rare enough that analyzing them to infer what investors are picking up can be a significant contribution in itself.


\section*{Vikas' Comments}

Vikas asks, what if livestreams attract capital to livestream funds from outside investors, as well as prompt (inside) investors of non-livestream funds to withdraw capital. In other words, capital inflows can be from different investors than those driving capital outflows. This is a fair point, which I don't think I have thought about.

But I'm still not sure why this would be an issue: a fund family might worry that it will shrink because (inside) investors of non-livestream funds will withdraw capital, so contemporaneously, it advertises livestream funds to cancel things out. This would be an observationally equivalent story, and I don't think we have to take a strong position.

Anywho, almost all of his comments relate to this issue.
\begin{itemize}
    \item Who's watching the livestreams? Is it possible to get even a rough sense of who's watching? Existing investors? New investors? Old investors? Young investors? I told him, I don't think this is something we have data to show.
    \item Why do fund families have to use social media livestreams instead of, e.g., counseling, etc.? Quintessentially, can you show livestreams are a more effective form of marketing than others? I think this is an unreasonably high bar.
    \item Which funds does a fund family market? He says, we make it sound like a fund family is marketing funds with high-performance. He thinks, if so, we should find that a fund family is marketing funds with the very best performance. If we find other characteristics that explain whether a fund family is marketing, it will be a different story. I think this is an unreasonably high bar: we are arguing that performance is a first-order characteristic for whether a fund is advertised or not, not necessarily that this is the only characteristic.
    \item He really wants to know the flow-weighted performance of funds that grow v. shrink around livestreams. Assuming that investors move from non-livestream funds to livestream funds, he thinks livestream funds should still beat non-livestream funds after the livestreaming event because of tax issues. I think he thinks there's some taxes when investors sell non-livestream funds to buy into livestream funds, so it should outperform. I don't think this is necessarily true, especially if we take his alternative story..
    \item He asks, what happens to the interaction of (i) whether a fund is livestreamed or not and (ii) whether a fund is high performer or not. Quintessentially, do non-livestream funds with high performance lose money simply because it was not livestreamed? Or do non-livestream funds with high performance also gain money because the livestreaming event awakes investors to examine options within the family. I think this is a fair point, and I think we should do it.
    \item Can you look at the effects of the livestream events on the next quarter rather than the contemporaneous quarter? He worries, the action may spill over to the next quarter, and we might be missing something there. I doubt we'll find any action there, but I think it's worth checking.

\end{itemize}

I forgot to add, he liked the story, but he thinks we really lack the data to confirm everything he's asking, and asked whether there's any chance we can get significantly more granular data for just one fund family. He says, if there's a way to get higher-frequency data and/or investor-level data for just one fund family would dramatically increase the potential, and he advises, that'd be what he'd focus on if he's working on the paper



\section*{The To Do List}

\begin{itemize}
    \item References to deHaan, Kacperczyk, and Song.
    \item What if double-cluster by family and quarter. This is how the literature tends to cluster when running flow regressions, although deHaan do otherwise. Of course, to the extent that our results are non-significant results, it's likely to make our results stronger anyways.
    
    The results are weaker (due to more conservative standard errors), but the key findings remain statistically significant and robust

    \item  \textcolor{red}{JK: It would be nice if we could compute the number of funds featured within a family-quarter / the total number of funds within a family-quarter and provide an average number here rather than come up with 20 funds example.}

    \item \textcolor{red}{By construction, 16\% of funds rank in the top five performers within their families (\textit{Fund Top 5 Return}). Why??}

    \item \textcolor{red}{JK: Isn't total fee = management fee + service fee + loads? If yes, is there a reason why we do not focus on loads alone (Table A6), just as we focus on management fee alone, and service fee alone? If total fee = management fee + service fee only, Table~\ref{tab:family_twfe_totalfee_reallocation} seems a bit redundant. Instead, I think it would make more sense to do fund-level analysis table like Table 3 for Table 2, where we include both \textit{Fund Livestream} and \textit{Family Livestream (Oth)} variables.}



\end{itemize}

Over the last few decades, the asset management industry has grown substantially in the U.S. and worldwide in terms of both the total assets under management (AUM) and the number of funds \citep{pastor2015scale,ici2025:ch1}. Given the proliferation of funds, a first-order friction in capital allocation that has been widely discussed in the literature is investor search costs \citep{hortacsu2004,huang2007participation,roussanov2021marketing}. This competition for investor flows has spurred fund families to engage aggressively in marketing: U.S. mutual funds, for example, spend roughly one-third of their fee revenue on marketing-related expenses, and such expenditures do attract more fund flows \citep{sirri1998costly,huang2007participation,gallaher2008advertising,christofferson2013what}. Despite the fact that marketing decisions are made at the fund family level, much of the existing empirical work evaluates the success of marketing efforts at the fund level. In particular, fund-level marketing ``success'' translates into fund-family-level asset growth only to the extent that inflows into marketed funds come from outside the family, rather than being reallocated from other funds within the same family. If marketing-induced inflows largely reflect within-family capital reallocation rather than new money, this would call for examining the motives behind mutual fund marketing rather than simply assuming that it is a growth strategy for fund families.




\textcolor{red}{Ringgenberg does some null result confirmation analysis to try to reassure it's not due to lack of statistical power, we might want to implement that going forward.}

\textcolor{red}{JK: Why is this interpretation correct? If we weigh by absolute flow magnitude, we would overweigh funds that received a lot of inflows AND funds that received a lot of outflows.}
Table~\ref{tab:family_twfe_churnweightedfee_reallocation} examines whether families strategically direct the largest flows toward higher-fee products by weighting fees by absolute flow magnitudes rather than assets. The results parallel the asset-weighted analysis: significant contemporaneous effects that disappear one quarter ahead. During livestreaming quarters, the most active flows are directed toward modestly higher-fee funds, but the lack of persistence indicates this is not a primary mechanism.

\end{document}
